% Options for packages loaded elsewhere
\PassOptionsToPackage{unicode}{hyperref}
\PassOptionsToPackage{hyphens}{url}
%
\documentclass[
  11pt,
]{article}
\usepackage{amsmath,amssymb}
\usepackage{iftex}
\ifPDFTeX
  \usepackage[T1]{fontenc}
  \usepackage[utf8]{inputenc}
  \usepackage{textcomp} % provide euro and other symbols
\else % if luatex or xetex
  \usepackage{unicode-math} % this also loads fontspec
  \defaultfontfeatures{Scale=MatchLowercase}
  \defaultfontfeatures[\rmfamily]{Ligatures=TeX,Scale=1}
\fi
\usepackage{lmodern}
\ifPDFTeX\else
  % xetex/luatex font selection
\fi
% Use upquote if available, for straight quotes in verbatim environments
\IfFileExists{upquote.sty}{\usepackage{upquote}}{}
\IfFileExists{microtype.sty}{% use microtype if available
  \usepackage[]{microtype}
  \UseMicrotypeSet[protrusion]{basicmath} % disable protrusion for tt fonts
}{}
\makeatletter
\@ifundefined{KOMAClassName}{% if non-KOMA class
  \IfFileExists{parskip.sty}{%
    \usepackage{parskip}
  }{% else
    \setlength{\parindent}{0pt}
    \setlength{\parskip}{6pt plus 2pt minus 1pt}}
}{% if KOMA class
  \KOMAoptions{parskip=half}}
\makeatother
\usepackage{xcolor}
\usepackage[margin=1in]{geometry}
\usepackage{graphicx}
\makeatletter
\def\maxwidth{\ifdim\Gin@nat@width>\linewidth\linewidth\else\Gin@nat@width\fi}
\def\maxheight{\ifdim\Gin@nat@height>\textheight\textheight\else\Gin@nat@height\fi}
\makeatother
% Scale images if necessary, so that they will not overflow the page
% margins by default, and it is still possible to overwrite the defaults
% using explicit options in \includegraphics[width, height, ...]{}
\setkeys{Gin}{width=\maxwidth,height=\maxheight,keepaspectratio}
% Set default figure placement to htbp
\makeatletter
\def\fps@figure{htbp}
\makeatother
\setlength{\emergencystretch}{3em} % prevent overfull lines
\providecommand{\tightlist}{%
  \setlength{\itemsep}{0pt}\setlength{\parskip}{0pt}}
\setcounter{secnumdepth}{5}
\ifLuaTeX
  \usepackage{selnolig}  % disable illegal ligatures
\fi
\usepackage{bookmark}
\IfFileExists{xurl.sty}{\usepackage{xurl}}{} % add URL line breaks if available
\urlstyle{same}
\hypersetup{
  pdftitle={Cantonal Engagement in the Prevention and Control of Vector-Borne Diseases in Switzerland and Liechtenstein},
  pdfauthor={Lilian Goepp, Nina Huber, Arlette Szelecsenyi, Julien Riou},
  hidelinks,
  pdfcreator={LaTeX via pandoc}}

\title{Cantonal Engagement in the Prevention and Control of Vector-Borne
Diseases in Switzerland and Liechtenstein}
\author{Lilian Goepp, Nina Huber, Arlette Szelecsenyi, Julien Riou}
\date{2025-10-29}

\begin{document}
\maketitle

{
\setcounter{tocdepth}{2}
\tableofcontents
}
\section{Introduction}\label{introduction}

This survey is part of the health project within the research program
``Decision Support for Dealing with Climate Change in Switzerland: a
cross-sectoral approach'' (NCCS-Impacts). The goal of NCCS-Impacts is to
develop practical climate-related services for the environment, the
economy, and society. The associated health project, jointly led by the
Federal Office of Public Health (FOPH) and the Federal Food Safety and
Veterinary Office (FSVO), focuses on the impacts of heat, mycotoxins,
and vector-borne diseases on both animal and human health.

Regarding vector-borne diseases, a comprehensive risk assessment of
current and future threats to humans and animals in Switzerland is being
carried out. As part of this risk assessment, the project documents (1)
the implementation of measures at the cantonal level in Switzerland and
in Liechtenstein and (2) the state of knowledge about vector-borne
diseases among the Bernese population. The national survey of
departments of health, veterinary services, and the environment aims to
gather information on engagement, specific measures, and collaboration
with other stakeholders.

\section{Methods}\label{methods}

To assess the engagement of cantonal and Liechtenstein authorities with
regard to vector-borne diseases, a cross-sectional survey was conducted
among stakeholders in departments responsible for health, veterinary
services, and environmental protection. In order to facilitate open
communication and increase willingness to participate, the responses
were not attributed to specific cantons (or Liechtenstein) or specific
departments outside of the general area of human health, animal health
and environment. Consequently, our analysis is limited in that it does
not identify which canton or department provided specific information,
nor does it allow for comparison between cantonal approaches.

The survey was implemented using the REDCap electronic data capture
platform (Research Electronic Data Capture, hosted by Unisanté) and
administered via an online questionnaire. The questionnaire was
distributed by email to targeted respondents within relevant cantonal
and Liechtenstein services, based on their role and responsibilities in
managing public health or environmental risks. Participation was
voluntary. Participants were able to save their answers and come back to
it later.

The online survey was open from April 25 to June 6, 2025. A reminder was
sent on May 16. The questionnaire included structured multiple-choice
and checkbox items, as well as a small number of open-ended questions,
covering topics such as awareness of vector-borne diseases, existing
preventive measures, interdepartmental collaboration, and perceived
needs or challenges. Respondents were ensured not to be identified.

The full questionnaire is available online:
\url{https://redcap.unisante.ch/surveys/?s=YKERXDNR4KAMWYP9}. Data were
anonymised, exported securely from REDCap and described using R (version
4.4.2). Open-text answers in multiple languages were summarized using
GPT-4-turbo (June 2025 version). More details are available in the
supplementary file.

\section{Results}\label{results}

\subsection{Participants}\label{participants}

We conducted a nationwide survey across all cantons, initially
contacting 88 general email addresses from the health, environment and
veterinary departments. Following feedback to include neobiota-related
stakeholders, we expanded the environmental sector contact list and sent
reminder emails to additional addresses, bringing the total to 153
contacts. Some recipients responded indicating their department was not
responsible for the survey topics, and others redirected us to
appropriate departments. Ultimately, we received 55 complete responses,
covering all 26 Swiss cantons and Liechtenstein. The number of answers
per canton varied between 1 and 3. Participants primarily belonged to
human health cantonal departments (n= 22), followed by the animal health
departments (n= 17) and the environment departments (n= 16).

\subsection{Involvement and resources}\label{involvement-and-resources}

Most participants (77\%) reported that their authority is currently
involved in the implementation of measures or activities related to
diseases transmitted by ticks and mosquitoes, such as public relations,
monitoring, or control (Supplementary table 1). This was especially true
in animal health departments (94\%). Among 12 of 55 respondents who
indicated no current involvement, the most frequently cited reasons
included a lack of financial or human resources (50\%) and a perception
that such diseases are outside the authority's responsibilities (42\%).
Regarding the existence of cantonal strategies or action plans, 49\% of
respondents were aware of a plan for mosquito-borne diseases and 11\%
for tick-borne diseases, while 33\% reported that no strategy exists in
their canton. Around half of the respondents reported lack of human
(49\%) and financial (55\%) resources

\subsection{Concrete measures and
activities}\label{concrete-measures-and-activities}

Reported concrete measures and activities in the field of vector-borne
diseases varied markedly between tick- and mosquito-borne diseases
(Supplementary table 2). The most common actions cited for ticks
included public information campaigns (22\%) and targeted awareness of
professionals (13\%), with rare reports of monitoring (4\% passive; 0\%
active) and vaccination advice (6\%). In contrast, involvement in
mosquito-related activities was broader and more diverse: 40\% of
respondents reported engaging in public relations and professional
awareness efforts, while active and passive mosquito monitoring were
cited by 36\% and 33\% of respondents, respectively. Nearly one-third
(29\%) mentioned verifying sightings at new locations, and around
15--18\% reported involvement in mosquito control on public or private
land. Information and training activities were mentioned by 20\% of
respondents. If the authority develops information, topics covered
included vaccination (65\%), protection against bites (53\%), and
elimination of breeding sites (47\%). Open-text responses further
highlighted involvement in national vaccination campaigns (e.g., against
tick-borne encephalitis for humans or bluetongue disease for ruminants
and camelids), interdepartmental coordination, media engagement, and
surveillance of animal diseases transmitted by ticks and mosquitoes as
defined in the Animal Disease Ordinance. Participants reported using a
wide range of channels to disseminate information, most commonly
official websites, emails, flyers, and direct communication with
veterinarians or animal owners. Additional methods included newsletters,
social media, articles in agricultural media, information sessions, and
training courses, often in collaboration with municipalities or other
local actors.

Respondents reported using a wide variety of guidelines, support
materials, and reference documents to inform the development and
implementation of measures against vector-borne diseases. The most
frequently cited sources included technical instructions and documents
from the Federal Food Safety and Veterinary Office (FSVO/OSAV/BLV), the
Federal Office for the Environment (FOEN/OFE/BAFU), and the Federal
Office of Public Health (FOPH/OFSP/BAG). Several participants also
referred to cantonal materials (particularly from Ticino) as well as
international resources such as those from the World Health Organization
(WHO), World Organization for Animal Health (WOAH), or cross-border
exchanges with France and Germany. Additionally, documents from academic
or technical partners like Scuola universitaria professionale della
Svizzera italiana (SUPSI) and Schweizerisches Tropen- und Public
Health-Institut (SwissTPH) were mentioned. A few respondents reported
using self-developed materials or relying on exchanges of expert
knowledge, while others indicated they did not use any specific
documents.

Respondents identified a range of knowledge gaps within their
authorities concerning both tick- and mosquito-borne diseases
(Supplementary table 3). For tick-borne diseases, the most commonly
cited areas of limited knowledge included disease management (20\%),
followed by detection and prevention (both 18\%). Open-text responses
often emphasized limited relevance or jurisdictional responsibility,
while some highlighted a lack of resources or the need for clearer
coordination and federally supported monitoring. For mosquito-borne
diseases, the main knowledge gaps were in disease and vector management
and control (36\%), as well as prevention (22\%) and detection (20\%).
Several respondents pointed to uncertainties around the appropriate use
of adulticides, sharing of responsibilities among stakeholders, and
long-term control strategies. Regarding actual control measures against
mosquitoes, over half of participants (53\%) reported that no active
mosquito control measures are in place. Among those that did report
activities, biological larvicides (25\%) and elimination of breeding
sites (20\%) were the most common, while use of adulticides (11\%) and
non-chemical larvicides (7\%) was less frequent. Open-text answers
revealed a mix of pilot studies, public awareness campaigns, and
experimental approaches like sterile insect techniques (SIT), often in
collaboration with academic or federal partners. Some respondents
indicated that control responsibilities lie with other agencies, such as
municipalities or environmental offices.

\subsection{Collaboration and
coordination}\label{collaboration-and-coordination}

Collaboration and coordination emerged as important themes in
participants' responses regarding vector-borne disease control
(Supplementary table 4). While 60\% of respondents reported active
networks or exchanges with researchers and specialists for
mosquito-borne diseases, only 20\% indicated such collaborations for
tick-borne diseases, and over a quarter (27\%) reported no active
exchange at all. Open-text responses highlighted needs for stronger
inter-cantonal coordination, improved knowledge transfer, and harmonized
federal guidelines, particularly for outbreak preparedness,
surveillance, and the responsible use of control tools like biocides or
adulticides. Several participants emphasized the importance of financial
and human resource support, centralized information platforms, and a
coordinated national monitoring strategy. Others cited challenges
related to fragmented responsibilities across levels of government or
lack of expertise and capacity. When asked about the integration of
tick- and mosquito-borne disease issues into national and cantonal
climate adaptation strategies, most respondents indicated limited
awareness or engagement: 65\% were unsure about national-level
integration, and 53\% did not know whether these issues were addressed
in their own canton's strategy.

\section{Conclusion}\label{conclusion}

This national survey provides a first detailed overview of how cantonal
authorities in Switzerland and Liechtenstein are engaged in the
monitoring, prevention, and control of diseases transmitted by ticks and
mosquitoes. While most respondents (77\%) reported some level of
involvement in related activities, findings reveal notable disparities
in resources, actions, and strategic alignment across cantons and
sectors.

Efforts appear more advanced for mosquito-borne diseases, with a broader
range of implemented activities such as public awareness campaigns,
professional training, and both active and passive monitoring. In
contrast, tick-related activities were reported far less frequently,
with more than half of respondents stating that their authority does not
currently engage in any tick-related measures. Access to dedicated human
or financial resources remains limited overall, particularly for
tick-borne disease interventions.

The survey highlights important knowledge gaps reported by cantonal
authorities, particularly in the areas of disease management and
pathogen detection. These gaps were reported more frequently for
mosquito-borne than tick-borne diseases, reflecting the growing public
health importance of invasive mosquito species. Despite these
challenges, several authorities reported innovative or pilot measures,
such as the use of biological larvicides and participation in sterile
insect technique studies.

Collaborative networks are more commonly established around
mosquito-related issues, with 60\% of respondents indicating active
exchanges with researchers and specialists in that domain, compared to
only 20\% for tick-related issues. Many participants called for stronger
coordination across cantons, improved knowledge transfer, and clearer
federal leadership, especially regarding outbreak preparedness,
standardized intervention strategies, and resource sharing.

Finally, the survey underscores a general lack of integration of
vector-borne diseases into climate adaptation strategies, both at the
national and cantonal levels. Most respondents were unaware of whether
these health issues are considered in existing climate strategies,
suggesting a need to strengthen the connection between climate and
vector-borne disease planning. Together, these findings point to the
importance of reinforcing institutional capacity, inter-cantonal
collaboration, and federal support in order to build a more coherent and
effective response to the rising threat of vector-borne diseases under
changing climatic conditions.

\end{document}
